\documentclass[11pt]{article}

\usepackage[margin=1in]{geometry}
\usepackage{amsmath}
\usepackage{amssymb}
\usepackage{booktabs}
\usepackage{hyperref}
\usepackage{array}

\title{Detailed Explanation of Model Architecture\\
\large Math 199 Crypto Pair-Spreads Project}
\author{Jack Lutz, Sammy Adham, and Frank Kronewitter}
\date{\today}

\begin{document}

\maketitle

\begin{center}
\fbox{\parbox{0.92\linewidth}{
\textbf{Draft note.} This document explains the architectures implemented in
the current repository. It is a technical draft for project documentation, not
a final methods section. If model code changes, this file should be checked
against \texttt{src/modeling.py}, \texttt{src/modeling\_big\_transformer.py},
\texttt{src/kalman\_hedge.py}, and \texttt{src/hmm\_filter.py}.
}}
\end{center}

\section{Prediction Task}

Each model predicts a three-class label for a pair spread over a 24-hour
horizon. Let \(s_t\) be the current spread and let \(s_{t+24}\) be the future
spread. The class label is based on whether the absolute spread moves closer
to zero, stays approximately unchanged, or moves farther away:
\[
y_t =
\begin{cases}
0, & |s_{t+24}| - |s_t| < -\tau \quad \text{(revert)},\\
1, & \left||s_{t+24}| - |s_t|\right| \leq \tau \quad \text{(persist)},\\
2, & |s_{t+24}| - |s_t| > \tau \quad \text{(diverge)}.
\end{cases}
\]
The fixed threshold is \(\tau=0.001\) in log-spread units unless the per-pair
label mode is enabled, in which case \(\tau\) is estimated from the pair's
training-period label volatility.

\section{Inputs}

The sequential models receive a 168-hour window:
\[
X_t \in \mathbb{R}^{168 \times 16}.
\]
The 16 per-hour feature channels are:
\[
\begin{array}{llll}
\text{spread}, & \text{spread\_z\_168h}, & \text{spread\_diff\_1h}, &
\text{spread\_diff\_4h},\\
\text{spread\_diff\_24h}, & \text{spread\_vol\_24h}, &
\text{spread\_vol\_168h}, & \text{volume\_ratio},\\
\text{volume\_z\_a\_168h}, & \text{volume\_z\_b\_168h}, &
\text{pair\_volume\_sum}, & \text{pair\_volume\_gmean},\\
\text{btc\_return\_24h}, & \text{rolling\_corr\_24h}, &
\text{realized\_vol\_a\_24h}, & \text{realized\_vol\_b\_24h}.
\end{array}
\]

The tabular model receives summary statistics computed from the same
168-hour window. These include latest, mean, minimum, and maximum spread
z-score; spread slope; latest and mean volume ratio; BTC 24-hour return;
rolling correlation; realized volatilities; time since zero crossing; and pair
volume summaries.

For learned models, features are standardized using training-window means and
standard deviations only:
\[
\tilde{x}_{i,j} = \frac{x_{i,j} - \mu_j}{\sigma_j}.
\]
The same training statistics are then applied to the test window.

\section{Model Summary}

\begin{table}[h]
\centering
\caption{Implemented model families.}
\begin{tabular}{llll}
\toprule
Model & Input & Trainable? & Output \\
\midrule
Persist baseline & none & no & class 1 \\
Majority baseline & labels only & no & training mode \\
Random stratified & labels only & no & sampled class \\
Z-score rule & latest spread z & no & class 0 or 1 \\
XGBoost / HGB & 14 tabular features & yes & 3-class probability \\
LSTM & \(168 \times 16\) sequence & yes & 3 logits \\
Small transformer & \(168 \times 16\) sequence & yes & 3 logits \\
Big transformer & \(168 \times 16\) sequence & yes & 3 logits \\
Kalman hedge & log-price pair & yes & dynamic residuals \\
Gaussian HMM & 3 spread features & yes & regime state \\
\bottomrule
\end{tabular}
\end{table}

\section{Non-Learned Baselines}

\subsection{Persist-Class Baseline}

The persist baseline always predicts class 1:
\[
\hat{y}_t = 1.
\]
In the trading convention, class 1 means no new spread trade. This baseline is
important because the class distribution can be imbalanced; any learned model
should improve on a no-action predictor.

\subsection{Majority-Class Baseline}

The majority baseline predicts the most common class in the training window:
\[
\hat{y}_t = \operatorname{mode}\{y_i : i \in \mathcal{T}_{\text{train}}\}.
\]
It uses no features. It tests whether a model is doing more than exploiting
the unconditional class distribution.

\subsection{Random-Stratified Baseline}

The random baseline samples predictions from the empirical training class
frequencies:
\[
\Pr(\hat{y}_t=k) =
\frac{\#\{i \in \mathcal{T}_{\text{train}} : y_i=k\}}
{\#\mathcal{T}_{\text{train}}}, \qquad k \in \{0,1,2\}.
\]
This gives a stochastic null model with the same marginal class distribution
as the training data.

\subsection{Z-Score Rule}

The z-score rule is the main classical trading baseline. It predicts
reversion when the absolute spread z-score is large:
\[
\hat{y}_t =
\begin{cases}
0, & |z_t| \geq 1.5,\\
1, & |z_t| < 1.5.
\end{cases}
\]
It never predicts class 2 directly. In the current results, this simple rule
has high per-trade win rate because it is selective.

\section{Tabular Booster}

\subsection{Feature Representation}

The booster uses the 14 tabular summary features from the 168-hour window:
\[
x_t \in \mathbb{R}^{14}.
\]
These features compress the recent spread path and market context into a
single row per sample. The row is standardized with a training-only
\texttt{StandardScaler} before fitting.

\subsection{XGBoost Architecture}

When \texttt{xgboost} is available, the model is an
\texttt{XGBClassifier} with:
\[
\begin{array}{ll}
\text{max depth} & 4,\\
\text{number of trees} & 300,\\
\text{learning rate} & 0.05,\\
\text{subsample} & 0.9,\\
\text{column sample by tree} & 0.9,\\
\text{objective} & \texttt{multi:softprob},\\
\text{evaluation metric} & \texttt{mlogloss}.
\end{array}
\]
Each tree partitions the tabular feature space. The ensemble adds trees
sequentially to reduce multiclass log loss:
\[
F_m(x) = F_{m-1}(x) + \eta f_m(x),
\]
where \(\eta=0.05\) is the learning rate and \(f_m\) is the \(m\)-th tree.
The final prediction is the class with the largest predicted probability:
\[
\hat{y}_t = \arg\max_{k \in \{0,1,2\}} p_k(x_t).
\]

\subsection{Fallback Histogram Gradient Booster}

If \texttt{xgboost} cannot be imported, the repository falls back to
\texttt{HistGradientBoostingClassifier} with 300 boosting iterations,
learning rate 0.05, and 31 maximum leaf nodes. This preserves the same
tabular-boosting role even when the external XGBoost dependency is unavailable.

\section{LSTM Classifier}

\subsection{Architecture}

The LSTM receives the full standardized sequence:
\[
X_t = (x_{t-167}, \ldots, x_t), \qquad x_i \in \mathbb{R}^{16}.
\]
The implemented architecture is:
\[
\begin{array}{ll}
\text{input size} & 16,\\
\text{hidden size} & 64,\\
\text{number of recurrent layers} & 2,\\
\text{dropout between LSTM layers} & 0.2,\\
\text{classifier head} & \mathbb{R}^{64} \rightarrow \mathbb{R}^{3}.
\end{array}
\]

At each time step, an LSTM cell updates a hidden state \(h_t\) and cell state
\(c_t\). In simplified notation,
\[
h_i, c_i = \operatorname{LSTMCell}(x_i, h_{i-1}, c_{i-1}).
\]
After processing all 168 hours, the model keeps the final hidden vector from
the last layer:
\[
h_{\text{last}} \in \mathbb{R}^{64}.
\]
The classifier head maps this vector to three logits:
\[
\ell_t = W h_{\text{last}} + b, \qquad \ell_t \in \mathbb{R}^{3}.
\]
The predicted class is \(\arg\max_k \ell_{t,k}\).

\subsection{Training}

The LSTM is trained with cross-entropy loss:
\[
\mathcal{L} = -\log
\frac{\exp(\ell_{t,y_t})}{\sum_{k=0}^{2}\exp(\ell_{t,k})}.
\]
The optimizer is Adam with learning rate \(10^{-3}\), weight decay
\(10^{-5}\), batch size 256, and the configured number of epochs in the
walk-forward script. In the headline deep walk-forward run, the script uses 4
epochs per split.

\section{Small Transformer Classifier}

\subsection{Architecture}

The small transformer is a compact encoder intended to be roughly comparable
in size to the LSTM. Its architecture is:
\[
\begin{array}{ll}
\text{input projection} & \mathbb{R}^{16} \rightarrow \mathbb{R}^{32},\\
\text{CLS token} & 1 \times 32 \text{ learned vector},\\
\text{encoder layers} & 2,\\
\text{attention heads} & 2,\\
\text{model dimension} & 32,\\
\text{feedforward dimension} & 64,\\
\text{dropout} & 0.2,\\
\text{classifier head} & \mathbb{R}^{32} \rightarrow \mathbb{R}^{3}.
\end{array}
\]

Each hourly feature vector is projected into a 32-dimensional embedding:
\[
e_i = W_{\text{in}}x_i + b_{\text{in}}.
\]
A learned classification token is prepended:
\[
Z_0 = [e_{\text{CLS}}, e_{t-167}, \ldots, e_t].
\]
The sequence is passed through two transformer encoder blocks. Each block
contains multi-head self-attention followed by a feedforward network. For a
single attention head,
\[
\operatorname{Attention}(Q,K,V)
=
\operatorname{softmax}\left(\frac{QK^\top}{\sqrt{d_k}}\right)V.
\]
The final CLS embedding is used for classification:
\[
\ell_t = W_{\text{head}} z_{\text{CLS}}^{(2)} + b_{\text{head}}.
\]

\subsection{Training}

The small transformer uses the same basic training loop as the LSTM:
cross-entropy loss, Adam, learning rate \(10^{-3}\), weight decay \(10^{-5}\),
batch size 256, and the configured walk-forward epoch count. In the current
post-audit headline run, this small transformer underperforms the LSTM and
booster; the later big-transformer experiment suggests that this is at least
partly a capacity and training-budget issue.

\section{Big Transformer Classifier}

\subsection{Architecture}

The big transformer is implemented separately to test whether the small
transformer result is caused by insufficient capacity. Its architecture is:
\[
\begin{array}{ll}
\text{input projection} & \mathbb{R}^{16} \rightarrow \mathbb{R}^{128},\\
\text{CLS token} & 1 \times 128 \text{ learned vector},\\
\text{positional encoding} & \text{fixed sinusoidal},\\
\text{encoder layers} & 4,\\
\text{attention heads} & 4,\\
\text{model dimension} & 128,\\
\text{feedforward dimension} & 512,\\
\text{activation} & \text{GELU},\\
\text{dropout} & 0.2,\\
\text{normalization} & \text{pre-norm encoder plus final LayerNorm},\\
\text{classifier head} & \mathbb{R}^{128} \rightarrow \mathbb{R}^{3}.
\end{array}
\]

The projected sequence is augmented with a CLS token and sinusoidal positions:
\[
Z_0 = [e_{\text{CLS}}, W_{\text{in}}x_{t-167}, \ldots, W_{\text{in}}x_t] + P,
\]
where \(P\) is a fixed sinusoidal positional encoding of length 169. The final
representation is the normalized CLS output after four encoder blocks:
\[
\ell_t = W_{\text{head}}\operatorname{LayerNorm}
\left(z_{\text{CLS}}^{(4)}\right) + b_{\text{head}}.
\]

\subsection{Training}

The big transformer uses a stronger training regime than the small transformer:
\[
\begin{array}{ll}
\text{optimizer} & \text{AdamW},\\
\text{learning rate} & 5 \times 10^{-4},\\
\text{weight decay} & 10^{-4},\\
\text{epochs} & 20,\\
\text{batch size} & 128,\\
\text{gradient clipping} & 1.0,\\
\text{scheduler} & \text{cosine decay with 1 warmup epoch},\\
\text{validation split} & \text{last 10\% of training data},\\
\text{early stopping} & \text{patience 5 on validation loss}.
\end{array}
\]
The validation split is time-ordered, not randomly shuffled, so it better
matches the forecasting setup.

\section{Kalman Dynamic Hedge Model}

\subsection{Purpose}

The Kalman model is not a three-class predictor. It is the dynamic hedge-ratio
model used to construct better spreads before feature engineering. It replaces
a static OLS hedge ratio with time-varying intercept and slope states.

\subsection{State-Space Model}

For log prices \(y_t=\log(P_{a,t})\) and \(x_t=\log(P_{b,t})\), the state is
\[
\theta_t =
\begin{bmatrix}
\alpha_t\\
\beta_t
\end{bmatrix}.
\]
The transition model is a random walk:
\[
\theta_t = \theta_{t-1} + w_t,
\qquad
w_t \sim \mathcal{N}(0,Q),
\]
with diagonal process covariance
\[
Q =
\begin{bmatrix}
q_\alpha & 0\\
0 & q_\beta
\end{bmatrix}.
\]
The observation model is:
\[
y_t =
\begin{bmatrix}
1 & x_t
\end{bmatrix}
\theta_t + v_t,
\qquad
v_t \sim \mathcal{N}(0,R).
\]
The Kalman residual is the one-step innovation:
\[
\epsilon_t = y_t -
\begin{bmatrix}
1 & x_t
\end{bmatrix}
\hat{\theta}_{t|t-1}.
\]
These residuals become the dynamic spread used by the downstream pipeline.

\subsection{Parameter Fitting}

Initial \(\alpha_0\) and \(\beta_0\) are seeded from OLS. The noise parameters
\(q_\alpha\), \(q_\beta\), and \(R\) are fit on the training slice by
minimizing the Kalman negative log likelihood:
\[
\mathcal{L}
=
\sum_t
\frac{1}{2}\left[
\log(2\pi S_t) + \frac{\epsilon_t^2}{S_t}
\right],
\]
where \(S_t\) is the innovation variance. Fitting uses L-BFGS-B in log
parameter space. In out-of-sample evaluation, the fitted parameters and final
training state are carried forward; the test window is not refit.

\section{Gaussian HMM Regime Filter}

\subsection{Purpose}

The HMM is a regime filter rather than a standalone predictor. It attempts to
detect whether a pair is in a mean-reverting regime or a diverging regime, then
suppresses trades outside the mean-reverting state.

\subsection{Inputs and Hidden States}

The HMM uses three spread-derived features:
\[
x_t =
\begin{bmatrix}
\text{spread\_z\_168h}_t\\
\text{spread\_diff\_1h}_t\\
\text{spread\_vol\_24h}_t
\end{bmatrix}.
\]
The default model has \(K=2\) hidden states. The hidden state follows a Markov
chain:
\[
z_t \sim \operatorname{Categorical}(A_{z_{t-1},:}),
\]
and observations are Gaussian conditional on the state:
\[
x_t \mid z_t=k \sim \mathcal{N}(\mu_k,\Sigma_k).
\]
The implementation uses full covariance matrices and expectation-maximization
through \texttt{hmmlearn}.

\subsection{Mean-Reverting State Identification}

After fitting, the model decodes states with Viterbi. The state with lower
standard deviation of \(\text{spread\_diff\_1h}\) is labeled
``mean-reverting.'' If there is a tie, the state with lower mean absolute
spread z-score is chosen. Predictions from other models are then filtered:
\[
\hat{y}^{\text{filtered}}_t =
\begin{cases}
\hat{y}_t, & z_t = z_{\text{MR}},\\
1, & z_t \neq z_{\text{MR}}.
\end{cases}
\]
Class 1 is the flat/no-trade class.

\section{Training and Walk-Forward Evaluation}

The main walk-forward driver trains models independently on each split:
\[
90 \text{ days train} \rightarrow 30 \text{ days test},
\]
then advances by 30 days. The long Kalman pipeline produces 27 splits. Each
split trains from scratch using only training-window samples and evaluates on
the held-out test window.

For deep models, the sequence tensor is indexed by sample id so every model
uses the same examples. For the tabular model, the same sample ids are mapped
to window-level summary features. This keeps model comparisons aligned across
the same pair-time observations.

\section{Prediction-to-Trade Convention}

Model outputs are class labels, not direct portfolio weights. The simple model
comparison metric translates classes into spread-direction signals:
\[
\text{signal}_t =
\begin{cases}
-\operatorname{sign}(s_t), & \hat{y}_t=0 \quad \text{(revert)},\\
0, & \hat{y}_t=1 \quad \text{(persist)},\\
\operatorname{sign}(s_t), & \hat{y}_t=2 \quad \text{(diverge)}.
\end{cases}
\]
This convention is used for comparing predictive signal quality. The separate
portfolio backtester adds position state, notional sizing, and transaction
costs.

\section{Interpretation}

The architecture stack is deliberately layered. The Kalman model defines a
dynamic residual spread. The feature builder converts that spread into
backward-looking sequence and tabular representations. Baseline, z-score,
booster, LSTM, and transformer models then predict whether the spread should
revert, persist, or diverge. Finally, the HMM can be applied as an optional
regime filter, and the backtester converts predictions into positions.

The current project results suggest that the strongest statistical advance is
the Kalman dynamic spread construction. The learned prediction models improve
some signal-quality metrics, but much of their useful information appears to
come from the latest spread z-score. This is why the architecture
documentation should be read together with the cost and selection-bias results
in the main paper draft.

\end{document}
